\documentclass[12pt, a4paper]{article}

\usepackage[english]{babel}
\usepackage[utf8x]{inputenc}
\usepackage{amsmath}
\usepackage{amssymb}
\usepackage{graphicx}
\usepackage{bbm}
\usepackage[colorinlistoftodos]{todonotes}
\usepackage{lmodern}
\usepackage{graphicx}

\usepackage{amssymb}
\usepackage{amsmath}
\usepackage{autobreak}
\usepackage{breqn}
\usepackage{hyperref}
\usepackage{ulem}
\usepackage{xcolor}
\usepackage{graphicx}  
\usepackage{subcaption} 
\usepackage{booktabs}
\newtheorem{proposition}{Proposition}

\title{Order Flow Imbalance Report}

\author{Ruming Liu}

\date{\today}

\begin{document}
\maketitle


\section{Tasks}

You must construct the following OFI features: 

\begin{itemize}
    \item Best-Level OFI 
    \item Multi-Level OFI
    \item Integrated OFI
    \item Cross-Asset OFI
\end{itemize}

\textcolor{blue}{I follow the methodology in the paper and use same time period setting. The OFI of each stock is generated in every single minute.}\\

\textcolor{blue}{(1) Best-Level OFI is generated under folder processd\_data/OFI\_0\_h\_i\_t/}\\

\textcolor{blue}{(2) Multi-Level OFI is generated under folder processd\_data/ofi\_m\_h\_i\_t/}\\

\textcolor{blue}{(3) Integrated OFI is generated under folder processd\_data/ofi\_I\_h\_i\_t/}\\

\textcolor{red}{(4) Cross-Asset OFI can be generated by using the same method in (1) and (3). But the data that I receieved is incomplete, it only includes the data of stock Apple Inc. I cannot generate cross-asset OFI without giving the cross-sectional information of other stocks.}


\section{Discussion}
\begin{itemize}
    \item Q1. What's the motivation behind measuring OFI at multiple depth levels of the order book?
    \item Q2. Why do the authors use Lasso regression rather than OLS for estimating cross-impact?
    \item Q3. Why is OFI considered a better predictor of short-term returns than trade volume?
    
    \item \textcolor{blue}{A1. Naturally, the order book at deeper depth (level 2 data) can capture more information instead of best bid and ask (level 1 data).
    Empirically, some algorithm trading strategies are layering, or hidden deeper in the order book. Muti-level OFI can capture such microstructure information. In the paper,
    the authors also find that muti-level OFI also capture the cross-impact of OFI of other stocks.}

    \item \textcolor{blue}{A2. LASSO decreases the model complexity and prevents over-fitting, which is a common problem in high-dimensional data. The sparse model
    is more acceptable in asset pricing.}
    \item \textcolor{blue}{A3. Trading volume is a technical indicator, which only includes the information of contemporaneity and history, from the hypothesis of efficient market, the past information should not be useful for predicting future returns.
    Also, volume is unsigned, which means it does not include the information of buy and sell. In contrast, OFI is a signed indicator, which can capture the imbalance of buy and sell orders, which has been proved 
    to be a useful factor in predicting future returns in the market microstructure literatures.}
\end{itemize}

\section{Reference}
\begin{itemize}
    \item[$\blacksquare$]Rama Cont, Mihai Cucuringu \& Chao Zhang (2023) Cross-impact of
    order flow imbalance in equity markets, Quantitative Finance, 23:10, 1373-1393, DOI: 10.1080/14697688.2023.2236159

\end{document}